% Section 4.3.1 - R(OC,sigma) shifted up by higher noise, convex mills ratio.
% GENERATED BY generate_figures.py - edit that file, not this one.
% sigma_1=0.9: OC_1=0.2758 (dR/dOC=+0.9064), stable (regular True, top True); no crossing from above
% sigma_2=2.25: OC_2=0.2448 (dR/dOC=+1.5439), Definition 3 conditions (regular True, top True)
\begin{figure}[htbp]
\centering
\pgfplotsset{ocshift/.style={
  width=12.0cm, height=8.4cm,
  xmin=0.13, xmax=0.33, ymin=0.381, ymax=0.673,
  axis x line=bottom, axis y line=left,
  axis line style={-{Stealth[length=6pt]},black}, clip=false,
  xlabel={$OC$}, ylabel={$\mathcal{R}(OC,\sigma)$},
  every axis x label/.style={at={(current axis.south east)},anchor=west,xshift=9pt},
  every axis y label/.style={at={(current axis.north west)},anchor=south,yshift=7pt},
  tick label style={font=\footnotesize}, label style={font=\footnotesize},
  title style={font=\footnotesize, yshift=2pt},
  tick style={black},
  xtick={0.1400,0.2448,0.2758}, xticklabels={,,},
  ytick={0.5700}, yticklabels={$\lambda^{R}$},
  tick label style={font=\small}, label style={font=\small},
}}

\begin{tikzpicture}
\begin{axis}[ocshift]
  \addplot[black,dashed,thick,samples=2,domain=0.13:0.33] {0.5700};
  \draw[gray,dashed] (axis cs:0.24476,0.381) -- (axis cs:0.24476,0.57000);
  \draw[gray,dashed] (axis cs:0.27580,0.381) -- (axis cs:0.27580,0.61688);
  \node[anchor=north,font=\small] at (axis cs:0.14000,0.3788) {$c$};
  \node[anchor=north,font=\small] at (axis cs:0.24476,0.3788) {$OC_2$};
  \node[anchor=north,font=\small] at (axis cs:0.27580,0.3788) {$OC_1$};
  \addplot[blue,very thick,smooth] coordinates {
(0.1442,0.4060) (0.1468,0.4097) (0.1495,0.4134) (0.1521,0.4171) (0.1547,0.4209)
(0.1574,0.4246) (0.1600,0.4283) (0.1627,0.4320) (0.1654,0.4357) (0.1681,0.4395)
(0.1708,0.4432) (0.1734,0.4469) (0.1761,0.4505) (0.1788,0.4542) (0.1815,0.4578)
(0.1842,0.4615) (0.1869,0.4651) (0.1895,0.4687) (0.1922,0.4722) (0.1948,0.4758)
(0.1975,0.4793) (0.2001,0.4828) (0.2027,0.4862) (0.2054,0.4896) (0.2080,0.4930)
(0.2105,0.4963) (0.2131,0.4996) (0.2157,0.5028) (0.2182,0.5060) (0.2207,0.5092)
(0.2232,0.5123) (0.2257,0.5154) (0.2281,0.5184) (0.2305,0.5213) (0.2329,0.5242)
(0.2353,0.5271) (0.2377,0.5298) (0.2400,0.5326) (0.2423,0.5352) (0.2446,0.5378)
(0.2468,0.5404) (0.2491,0.5429) (0.2512,0.5453) (0.2534,0.5476) (0.2555,0.5499)
(0.2576,0.5522) (0.2597,0.5543) (0.2617,0.5564) (0.2638,0.5585) (0.2657,0.5605)
(0.2677,0.5624) (0.2696,0.5642) (0.2715,0.5660) (0.2733,0.5677) (0.2752,0.5694)
(0.2770,0.5710) (0.2787,0.5726) (0.2804,0.5741) (0.2821,0.5756) (0.2838,0.5770)
(0.2854,0.5783) (0.2870,0.5797) (0.2886,0.5809) (0.2901,0.5822) (0.2916,0.5834)
(0.2931,0.5845) (0.2946,0.5857) (0.2960,0.5868) (0.2974,0.5878) (0.2987,0.5889)};
  \addplot[blue,very thick,dash pattern=on 6pt off 3pt,smooth] coordinates {
(0.1442,0.4062) (0.1468,0.4102) (0.1495,0.4143) (0.1521,0.4185) (0.1547,0.4228)
(0.1574,0.4271) (0.1600,0.4315) (0.1627,0.4360) (0.1654,0.4404) (0.1681,0.4449)
(0.1708,0.4495) (0.1734,0.4540) (0.1761,0.4585) (0.1788,0.4630) (0.1815,0.4675)
(0.1842,0.4720) (0.1869,0.4765) (0.1895,0.4810) (0.1922,0.4854) (0.1948,0.4899)
(0.1975,0.4942) (0.2001,0.4986) (0.2027,0.5029) (0.2054,0.5072) (0.2080,0.5115)
(0.2105,0.5157) (0.2131,0.5199) (0.2157,0.5240) (0.2182,0.5281) (0.2207,0.5321)
(0.2232,0.5361) (0.2257,0.5400) (0.2281,0.5439) (0.2305,0.5478) (0.2329,0.5516)
(0.2353,0.5553) (0.2377,0.5590) (0.2400,0.5626) (0.2423,0.5662) (0.2446,0.5697)
(0.2468,0.5732) (0.2491,0.5766) (0.2512,0.5799) (0.2534,0.5832) (0.2555,0.5865)
(0.2576,0.5897) (0.2597,0.5928) (0.2617,0.5959) (0.2638,0.5989) (0.2657,0.6019)
(0.2677,0.6048) (0.2696,0.6076) (0.2715,0.6104) (0.2733,0.6132) (0.2752,0.6159)
(0.2770,0.6186) (0.2787,0.6212) (0.2804,0.6238) (0.2821,0.6264) (0.2838,0.6289)
(0.2854,0.6314) (0.2870,0.6339) (0.2886,0.6363) (0.2901,0.6388) (0.2916,0.6412)
(0.2931,0.6436) (0.2946,0.6460) (0.2960,0.6484) (0.2974,0.6508) (0.2987,0.6532)};
  % the dots carry a white rim so they stand out against the curves
  \filldraw[fill=darkgreen,draw=white,line width=0.8pt] (axis cs:0.27580,0.57000) circle (3.2pt);
  % the new stable crossing
  \filldraw[fill=red,draw=white,line width=0.8pt] (axis cs:0.24476,0.57000) circle (3.2pt);
  % where the guide at OC_1 meets the new curve
  \filldraw[fill=orange,draw=white,line width=0.8pt] (axis cs:0.27580,0.61688) circle (3.2pt);
  \node[blue,anchor=west,font=\small] at (axis cs:0.3019,0.6532) {$\mathcal{R}(OC,\sigma_2)$};
  \node[blue,anchor=west,font=\small] at (axis cs:0.3019,0.5889) {$\mathcal{R}(OC,\sigma_1)$};
\end{axis}
\end{tikzpicture}
\caption{}
% Caption kept here, deliberately not typeset - the figure carries a
% number only.  To restore it, move the line below into \caption{}.
% The regular acceptance rate curve before (solid) and after (dashed) a rise in noise from $\sigma_1$ to $\sigma_2$. The mills ratio is convex, so the curve shifts upwards. It crosses $\lambda^R$ from below at $OC_1$, a stable equilibrium in the sense of Definition \ref{def:equilibrium}; after the shift, the crossing from below moves to the lower opportunity cost $OC_2$, which is also stable. $\sigma_1=0.9$, $\sigma_2=2.25$; $F$ standard logistic, $\sigma=0.9$; $G=\mathrm{Beta}(2,2)$ on $[-2.5,2.5]$; $\lambda^T=0.3$, $v^R/v^T=0.35$, $\lambda^R=0.57$, $c/v^R=0.4$.
\label{fig:oc_noise_shift}
\end{figure}
